\renewcommand{\Anno}{1999}

\renewcommand{\NumeroProblema}{1}

\renewcommand{\Testo}{Dato un numero reale $\alpha$, indichiamo con
$[\alpha]$ la sua parte intera, definita come il più grande intero minore od uguale ad $\alpha$.
	
Determinare tutte le soluzioni reali dell'equazione
$$x[x[x]]=84.$$}

\renewcommand{\Risposta}{L'equazione ha l'unica soluzione $x=14/3$.}

\renewcommand{\Aiutino}{Dimostare che $x$ deve stare tra 4 e 5, da cui
$[x]=4$.  L'equazione si riduce pertanto a $x[4x]=84$.  A questo punto $[4x]$ può assumere solo determinati valori, dai quali si ricava immediatamente $x$.  Non resta che controllare quali dei valori di $x$
trovati sono effettivamente soluzioni.}

\renewcommand{\SoluzioneUfficiale}{
}

\renewcommand{\SoluzioniAlternative}{
}

\renewcommand{\Osservazioni}{
}

\renewcommand{\Commenti}{
}
% Aggiunta redazionale: la soluzione preesistente non viene modificata.
\providecommand{\SoluzioneRedazionale}{}
\renewcommand{\SoluzioneRedazionale}{
% SSSUP-BOZZA-AI: SoluzioneRedazionale; autore=ChatGPT; data=2026-11-08; stato=DA-VERIFICARE
% SSSUP-PROVENIENZA: AI: elaborazione nuova dal testo
Interpretiamo la scrittura come $x\lfloor x\lfloor x\rfloor\rfloor=84$. Se $x<0$, allora $x\lfloor x\rfloor>0$ e quindi il membro sinistro è non positivo. Se $0\leq x<4$, esso è al più $x^3<64$; se $x\geq5$, è almeno $125$. Ne segue $4\leq x<5$ e dunque $\lfloor x\rfloor=4$.

Posto $k=\lfloor4x\rfloor$, abbiamo $k\in\{16,17,18,19\}$ e $x=84/k$. La condizione che definisce $k$ diventa
\[
k\leq\frac{336}{k}<k+1,
\qquad\text{ossia}\qquad k^2\leq336<k(k+1).
\]
Tra i quattro candidati solo $k=18$ soddisfa entrambe le disuguaglianze. Si ottiene $x=84/18=14/3$; sostituendo, $\lfloor x\rfloor=4$ e $\lfloor4x\rfloor=18$, per cui questa è effettivamente una soluzione e l'unica.
}
